\begingroup
\setlength{\abovedisplayskip}{6pt}
\setlength{\belowdisplayskip}{6pt}
\setlength{\abovedisplayshortskip}{3pt}
\setlength{\belowdisplayshortskip}{4pt}

% Printed 102; principal PDF 109; Google PDF 113.
\par\clearpage
\noindent\makebox[12mm][l]{\fontsize{13}{16}\selectfont 102}\makebox[\dimexpr\linewidth-24mm\relax][c]{\scriptsize\textbf{CORRESPONDENCE.\quad VARIATION OF THE ALTITUDE OF THE POLE WITH HEIGHT.}}\makebox[12mm][r]{}\par\medskip

\[
\theta=\frac{g'-g}{g}=\tfrac54\frac hg
\]
\noindent holds (\textsc{Newton}'s flattening).

In this new, more general form, if, with \textsc{Sabine}, I set
\[
\frac{g'-g}{g}=\frac1{192{,}7}
\]
\noindent we obtain for the altitude of the pole at \(a\)
\[
=\varphi+1070''\,.\,\frac{s}{a_0}\sin 2\varphi.
\]

Thus, for the highest mountains in Silesia, only about \(\tfrac14\) of a second. I must
frankly confess to you, however, that I can regard the entire investigation only
as a theoretical curiosity, wholly devoid of practical significance.
It would have such significance only if, on the smooth surface of the Earth
\(DAE\), a thin, tall column \(Aa\) were erected, at the top
and foot of which the altitude of the pole could be observed. In reality, when
\(a\) is, say, on a high mountain, the point \(A\) is in the first place
quite inaccessible. Even if one could reach it and could
precisely determine by measurement the difference in the direction of the plumb line at \(A\) and \(a\),
one would still have no grounds for regarding the above formula
as representing this difference, since the attractions of the parts
of the body of the Earth lying above the level of \(A\) will produce much larger differences
in the final results for gravity
at \(A\) and \(a\), differences that cannot be subjected to any calculation.

I have already expressed myself in general terms on this subject in my work
of 1828 on the difference in latitude between Göttingen and Altona, p.~73 [cf.\
this volume, p.~49], and I know of nothing better to say
about it now.

\begin{center}\rule{.18\textwidth}{.4pt}\end{center}
{\small
\begin{center}\textbf{REMARKS.}\end{center}
The preceding letter, of which a copy is held in the Gauss Archive, is reprinted in
»Protocollen der Verhandlungen der permanenten Commission der europäischen Gradmessung vom 23. bis
29. September 1869 in Florenz«. (Printed as a manuscript.) pp.~30---32.
\par\hfill\textsc{Krüger, Börsch.}\hspace{.06\textwidth}\par
}
\begin{center}\rule{.24\textwidth}{.4pt}\end{center}

